\section{System Overview}
\label{sec:SystemOverview}

Durga has three centers:
\begin{itemize}
\item a BPMN parsing and generation pipeline,
\item a generated Kafka-oriented process runtime,
\item a monitoring and reporting path built around canonical lifecycle events.
\end{itemize}

\subsection{Repository-level components}

\begin{description}
\item[Scaffolder] Parses a BPMN model and renders a generated project (entry point:
\texttt{org.gautelis.durga.tools.BpmnScaffolder}).
\item[Generation support] Model extraction and template coordination: \texttt{BpmnModelSupport},
\texttt{TaskRoutingGenerator}, \texttt{EventTemplateGenerator},
\texttt{SubProcessTemplateGenerator}, \texttt{GeneratedProjectSupport}.
\item[Templates] StringTemplate groups under \texttt{src/main/resources/templates/} define
generated Java classes, scripts, YAML, Maven build files and shared runtime helpers.
\item[Monitoring runtime] A Kafka Streams topology consumes canonical lifecycle events and
materializes state, counts, latency summaries and stuck-instance views.
\end{description}

\subsection{Generated runtime shape}

Each generated project contains:
\begin{itemize}
\item worker and routing services,
\item a starter,
\item a process orchestrator,
\item helpers for user/manual task completion and failure testing,
\item generated Kafka topic configuration,
\item sample payloads and runnable scenario scripts.
\end{itemize}

Topics and handlers are deliberately visible so that BPMN concepts are projected concretely
onto Kafka channels and lifecycle events.

\begin{figure}[!ht]
    \centering
    \begin{tikzpicture}[node distance=16mm,>=latex]
        \tikzstyle{box}=[draw=iptoMuted, rounded corners=2pt, line width=0.8pt, align=center, minimum width=34mm, minimum height=11mm, fill=white]
        \tikzstyle{accentbox}=[box, fill=iptoAccent!12]
        \node[accentbox] (bpmn) {BPMN model\\input};
        \node[box, right=22mm of bpmn] (scaffolder) {Durga scaffolder\\Camunda + ST4};
        \node[box, right=22mm of scaffolder] (generated) {Generated process\\project};

        \node[box, below=16mm of scaffolder] (events) {Canonical\\process-events};
        \node[box, right=22mm of events] (monitoring) {Kafka Streams\\monitoring};
        \node[box, left=22mm of events] (topics) {Execution topics\\and handlers};

        \draw[->, thick] (bpmn) -- (scaffolder);
        \draw[->, thick] (scaffolder) -- (generated);
        \draw[->, thick] (generated) -- (topics);
        \draw[->, thick] (generated) -- (events);
        \draw[->, thick] (events) -- (monitoring);
    \end{tikzpicture}
    \caption{High-level flow from BPMN model to generated runtime and monitoring projections}
    \label{fig:DurgaOverview}
\end{figure}

\subsection{Canonical event model}

The shared canonical stream on \texttt{process-events} is the observability backbone for
latest state per instance, counts by current state, latency summaries, stuck-instance
detection, dashboards and query endpoints. Execution topics may be useful for routing
visibility, but the monitoring model depends on explicit projections rather than raw topic
inspection.
